% Source-aware verification context for the attached P21.3 subcase only.
\section*{Scope and load-bearing inputs}

The target is a restricted continuous-pair P21.3 assertion on an open inverse
image in a split toric variety.  Unlike the prior dense-torus subcase, its
tropical open set may meet toric boundary strata.  It does not claim the
general P21.3 predicate, noninvariant potentials, nonsplit tori, arbitrary
toroidal charts, or an arbitrary analytic domain.

The proof uses the following precisely delimited inputs.
\begin{enumerate}
\item BGJKK25 equation~(5.5) identifies continuous functions on every open
subset $O\subseteq N_\Sigma$ with $S$-invariant continuous functions on its
analytic inverse image.
\item The test-form calculation in the proof of BGJKK25
Proposition~5.3.10(iv) gives the pushforward identity for the current of a
continuous invariant function.  Its displayed calculation uses factorization,
continuity, canonical-skeleton support of pulled-back top-degree test forms,
and pushforward/differential commutation; it does not use psh positivity.
\item A basic neighborhood of a point in the stratum $N(\sigma)$ has dense
part $\Omega+\sigma$ for open bounded convex $\Omega$.  It contains every
compactified ray in a direction $w\in\sigma$.  This is the neighborhood model
in BGJKK25, proof of Theorem~3.1.4.  On its dense convex part, zero
$\dprimeddprime$ current means an affine continuous representative.  Finite
continuity at all ray endpoints forces its linear functional to vanish on
$\sigma$.
\item In BGJKK25 \S3.3, smooth functions on a tropical toric partial
compactification are exactly the locally smooth functions constant towards
the boundary.  The preceding slope calculation therefore makes the descended
affine function smooth.  BGJKK25 Proposition~2.3.4 pulls smooth Lagerberg
forms back to smooth CLD forms and commutes with both differentials.
\item Finally, the proof itself uses the graded Leibniz rule.  For a smooth
function $f$ and a $\dprime$- and $\ddprime$-closed current $T$,
$\dprimeddprime(fT)=\dprime(\ddprime f\wedge T)
=\dprimeddprime f\wedge T$.  Hence this current vanishes when the smooth
Hessian of $f$ vanishes.
\end{enumerate}

Please reject the candidate if: the pushforward calculation has an omitted
psh or support hypothesis; a basic boundary neighborhood does not supply the
claimed compactified rays; finite continuity does not force the normal affine
slopes to vanish; tropical smoothness fails to pull back to CLD smoothness;
or the displayed graded Leibniz calculation has a sign or bidegree defect.
